% =====================================================================
%  II. LITERATURE SURVEY
%  Every factual claim about prior work below is taken from
%  notes/prior-art-register.md, which records what was verified and from
%  which source. Do not add a claim about a paper that is not in the
%  register without reading the paper first.
% =====================================================================
\section{Literature Survey}
\label{sec:literature}

The problem sits where three bodies of work meet: routing and cascading
across language models, control of reasoning effort at inference time, and
measurement of the energy that inference consumes. We review each, then
the small set of papers that combine two of them, and close with the gap
this study addresses.

\subsection{Routing and Cascading Across Models}
\label{sec:lit-routing}

The word \emph{routing} has two meanings in this area, and they should be
kept apart. Mixture-of-experts architectures route tokens between expert
sub-networks inside a single model~\cite{fedus2022switch}. This paper is
concerned with the other meaning: routing whole queries between separate
models, each of which could answer on its own. Moslem and Kelleher survey
that literature and organise it into routing, which commits to one model
before generation, and cascading, which tries models in
sequence~\cite{moslem2026routingsurvey}.

\emph{Cascades.} FrugalGPT sends a query through a sequence of models of
increasing price and stops when a learned scorer judges the answer
acceptable; its authors report matching the best individual model with up
to 98\% lower cost~\cite{chen2023frugalgpt}. AutoMix replaces the external
scorer with few-shot self-verification by the small model and decides
whether to escalate with a partially observable Markov decision process,
reporting cost reductions above 50\%~\cite{aggarwal2024automix}. AutoMix is
the closest published idea to the question asked here, because it treats
escalation as a sequential decision. Its decision state, however, is the
model's confidence in its own answer, not the cost already spent on failed
attempts, and it is an automated policy rather than a model of what a user
does. Both cascades share a structural cost: the cheap model must generate
before the system can decide, so every escalated query pays for at least one
discarded answer.

\emph{Predictive routers.} Hybrid LLM trains a router that estimates, from
the query alone, whether a small model's answer will be close in quality to
a large model's, and exposes a threshold that trades quality for cost at
test time, reporting 40\% fewer calls to the large model with no loss in
response quality~\cite{ding2024hybridllm}. RouteLLM learns a
strong-versus-weak decision from human preference data and reports more than
a twofold cost reduction~\cite{ong2024routellm}. RouterDC trains a query
encoder with dual contrastive losses so that queries sit near the models
that answer them well~\cite{chen2024routerdc}. UniRoute addresses the
practical problem that the model pool changes, routing to models unseen
during training and giving an excess-risk bound for the
procedure~\cite{jitkrittum2025universal}. Arch-Router instead aligns a
1.5B-parameter router with user-defined domain and action preferences so
that models can be added without retraining~\cite{tran2025archrouter}, and
Router-R1 moves to multi-round routing, training a model with reinforcement
learning to interleave its own deliberation with calls to other models and
to aggregate their outputs across rounds~\cite{zhang2025routerr1}.

\emph{Benchmarks.} Three benchmarks make routers comparable. RouterBench
provides a 36{,}497-prompt zero-shot split with outputs and per-inference
costs for 11 models, and oracle labels for the best model per
prompt~\cite{hu2024routerbench}. RouterEval scales the candidate pool to
more than 8{,}500 models and 200 million performance
records~\cite{huang2025routereval}. LLMRouterBench collects more than
400{,}000 instances over 21 datasets and 33 models, with code and data
released publicly~\cite{li2026llmrouterbench}.

Across this body of work, cost is denominated in API price, model calls or
tokens, and the comparison points are static single-model policies, random
assignment, an oracle, or another automated router or cascade.

\subsection{Reasoning Effort and Test-Time Compute}
\label{sec:lit-reasoning}

Chain-of-thought prompting showed that asking a model to write intermediate
steps improves multi-step reasoning~\cite{wei2022cot}, and self-consistency
improved it further by sampling several reasoning paths and taking a
majority vote over their answers~\cite{wang2023selfconsistency}. Both buy
accuracy with generated tokens. Reinforcement learning has since been used
to train models that produce long reasoning traces without being
prompted to~\cite{deepseek2025r1}, and simple interventions at decoding time
can lengthen or cut those traces: s1 introduced budget forcing, which ends
the thinking phase early or suppresses its end to extend
it~\cite{muennighoff2025s1}. Qwen3 builds the control into the model,
serving thinking and non-thinking behaviour from one set of weights with a
user-settable thinking budget~\cite{qwen3}.

Long traces are not always useful. Sui \emph{et al.} survey methods for
efficient reasoning and document the tendency of reasoning models to spend
many tokens on easy questions~\cite{sui2025stopoverthinking}. Alomrani
\emph{et al.} organise adaptive test-time compute into two tiers:
controllability under a fixed budget, and adaptiveness, where the budget
itself depends on the input~\cite{alomrani2025reasoningbudget}. Selecting a
reasoning mode per prompt is an instance of the second tier.

\subsection{Routing Over Reasoning Mode}
\label{sec:lit-mode-routing}

A recent group of papers routes over whether and how much to think, usually
within a single model. SynapseRoute trains a logistic-regression classifier
on query embeddings to send medical questions to thinking or non-thinking
mode of one dual-state model, reporting 36.8\% lower inference time, 39.66\%
fewer tokens and that about 58\% of the questions need no reasoning
mode~\cite{zhang2025synapseroute}. Think When Needed attaches a
lightweight router head that chooses Think or Non-Think before generation
for ranking tasks~\cite{guo2026thinkwhenneeded}. DART is training-free: it
generates two cheap non-thinking drafts, reads problem difficulty from their
agreement, and sets a per-query thinking budget accordingly, cutting
thinking tokens by 32--73\%~\cite{lee2026dart}. Learning When to Think moves
the decision inside the model, which emits NoThink, Short or Long as its
first token after training with GRPO~\cite{kassenaar2026learningwhentothink}.
HRBench compares prompt-based, external-routing and speculative switching
strategies across four training regimes, six models and five benchmarks, and
concludes that the preferred strategy depends on model scale and task
domain~\cite{ning2026hrbench}.

Route to Reason is the closest prior work to the grid studied here. It
selects jointly among seven open-source models and four reasoning strategies
under a budget constraint and reports more than 60\% fewer
tokens~\cite{pan2026routetoreason}. Its existence means that joint selection
of model and reasoning effort is not, by itself, a new idea, and we do not
claim it as one.

All of these papers depend on an assumption that is rarely tested: that a
model's thinking and non-thinking modes are genuinely different operating
points. Wang \emph{et al.} examined hybrid-thinking models directly and found
only partial separation, with reasoning behaviour leaking into non-thinking
outputs; their improved training recipe reduced non-thinking output length on
MATH500 from 1{,}085 to 585 tokens, a large gain that still falls short of
clean separation~\cite{wang2025demystifying}. A routing study whose grid
includes a reasoning-mode axis inherits this risk, which is why we test
separation on our own checkpoints before labelling
(Section~\ref{sec:gate-mode}).

\subsection{Energy and Carbon of Language-Model Inference}
\label{sec:lit-energy}

Strubell \emph{et al.} brought the energy and financial cost of training NLP
models to the field's attention~\cite{strubell2019energy}, and Schwartz
\emph{et al.} argued that efficiency should be reported alongside accuracy
as a first-class result~\cite{schwartz2020greenai}. Lacoste \emph{et al.}
provided a calculator for converting compute into emissions using regional
grid intensity~\cite{lacoste2019quantifying}, and Patterson \emph{et al.}
showed how strongly the emissions of large training runs depend on
hardware, data-centre efficiency and location~\cite{patterson2021carbon}.
For inference, Luccioni \emph{et al.} measured energy and carbon per
1{,}000 queries across task types and found wide variation between
them~\cite{luccioni2024powerhungry}.

Work on constrained hardware is closer to our setting. Camel reduces energy
on an NVIDIA Jetson AGX Orin by tuning GPU frequency and batch size for a
single model~\cite{xu2025camel}, and serving systems such as vLLM raise
throughput by managing key--value cache memory more
efficiently~\cite{kwon2023vllm}. Both are orthogonal to routing: they make
each configuration cheaper to run but do not decide which configuration to
run.

\subsection{Energy-Aware and Latency-Aware Routers}
\label{sec:lit-energy-routing}

Three routers treat energy or latency as part of the objective, and we read
each in full because they bear most directly on our claims. GreenServ routes
queries among 16 open models using task type, semantic cluster and text
complexity, measures GPU energy on an A100 with the Zeus library, and reports
a 31\% cumulative energy reduction with 22\% higher accuracy than random
routing~\cite{ziller2026greenserv}. It characterises its router's overhead
at about 7.77\,ms per query but does not subtract the router's energy from
the saving. EcoThink sends each query either to a retrieval path or to a
reasoning path using a distilled DistilBERT classifier, and reports a 40.4\%
average energy reduction across nine benchmarks and carbon of
1.32\,gCO$_2$e per query against 4.52 for GPT-4o~\cite{li2026ecothink}. Its
energy comes from CodeCarbon sampling combined with an average utilisation
of thermal design power, so it is partly estimated, and the router's own
cost is not charged. Yang \emph{et al.} route queries among edge experts on
RTX~4090 GPUs using two DistilBERT-based predictors and an actor--critic
policy, and optimise quality of service and latency without measuring
energy~\cite{yang2025qosrouting}.

Two observations follow. First, energy-aware routing is established, and so
is reporting carbon for it; neither is a contribution of this paper. Second,
these routers select a model or a pipeline. None of them routes over
reasoning mode inside a hybrid model, and none subtracts the router's own
measured energy from the energy it saves.

\subsection{Evaluation With Language-Model Judges}
\label{sec:lit-judges}

Open-ended prompts cannot be graded by string match, so routing studies that
include them rely on model judges. Zheng \emph{et al.} found that a strong
judge agreed with human preferences at a rate comparable to agreement between
humans, while also documenting position, verbosity and self-enhancement
biases~\cite{zheng2023judging}. The biases matter here because a judge drawn
from the same family as the routed models could favour them, which motivates
the out-of-family judge and the two-sided validation in
Section~\ref{sec:judge}.

\subsection{Positioning}
\label{sec:lit-gap}

Table~\ref{tab:prior-work} summarises the systems closest to this study.
Reading them together, four questions remain open in the literature we
located.

\begin{enumerate}
  \item \emph{The escalation tax is unmeasured.} Routers and cascades are
        compared with static, random, oracle or automated baselines. None
        formalises manual escalation, in which each failed attempt is billed
        in full and the prompt is resent to the next configuration.
  \item \emph{Router energy is not netted out.} Energy-aware routers report
        gross savings. We found no study that subtracts the selector's own
        measured energy, or that asks at what point a language-model router
        stops paying for itself against a linear classifier.
  \item \emph{The reasoning-mode axis has no measured energy cost.} Routers
        over reasoning mode report tokens and latency; routers that measure
        energy do not route over reasoning mode.
  \item \emph{Energy is not normalised by success.} Energy per query and per
        token are standard. Energy per \emph{solved} query, which charges a
        cheap failure for the retry it causes, was not found.
\end{enumerate}

\noindent The study also characterises the grid on an 8\,GB laptop GPU. That
is a useful regime to report, since the energy-aware routers above ran on
A100 and RTX~4090 hardware, but we treat it as context rather than as a
contribution in its own right.

% Table cells are drawn from notes/prior-art-register.md. Every row now
% agrees with the prose above, which states each system's own reported
% figures. Two rows were checked against the full texts because the prose
% did not pin the metric down: Router-R1's cost reward is output tokens
% weighted by a function of parameter count, and the QoS router's headline
% is a composite of quality and latency compliance rather than latency
% alone. The four 2023--2024 rows are long-established. Rows still worth a
% PDF pass before submission: SynapseRoute, DART, Learning When to Think
% and Route to Reason, whose decision points come from their abstracts.
\begin{table*}[!t]
\centering
\caption{Systems closest to this study. \emph{Axis} is what the router
selects; \emph{decision point} is when it commits; \emph{efficiency
reported} is the cost measure used in the headline result.}
\label{tab:prior-work}
\renewcommand{\arraystretch}{1.15}
\begin{tabularx}{\textwidth}{@{}l l >{\raggedright\arraybackslash}X >{\raggedright\arraybackslash}X@{}}
\toprule
System & Axis & Decision point & Efficiency reported \\
\midrule
FrugalGPT~\cite{chen2023frugalgpt} & Model & Cascade, scorer after each answer & API cost \\
AutoMix~\cite{aggarwal2024automix} & Model & Cascade, self-verification & Cost \\
Hybrid LLM~\cite{ding2024hybridllm} & Model (small/large) & Before generation & Large-model calls \\
RouteLLM~\cite{ong2024routellm} & Model (strong/weak) & Before generation & Cost \\
Router-R1~\cite{zhang2025routerr1} & Model & Multi-round & Parameter-weighted output tokens \\
SynapseRoute~\cite{zhang2025synapseroute} & Reasoning mode & Before generation & Tokens, time \\
DART~\cite{lee2026dart} & Thinking budget & After two short drafts & Tokens, latency \\
Learning When to Think~\cite{kassenaar2026learningwhentothink} & Thinking level & In-model, first token & Response length \\
Route to Reason~\cite{pan2026routetoreason} & Model and strategy & Before generation & Tokens \\
GreenServ~\cite{ziller2026greenserv} & Model & Before generation & Measured energy (gross) \\
EcoThink~\cite{li2026ecothink} & Retrieval or reasoning path & Before generation & Partly estimated energy, carbon (gross) \\
QoS edge routing~\cite{yang2025qosrouting} & Edge expert & Before generation & QoS score under a latency bound \\
\midrule
This study & Model size and reasoning mode & Before generation & Tokens, latency, measured energy (net of router) \\
\bottomrule
\end{tabularx}
\end{table*}
